\section{Operational context}

Situated in the municipality of Laterza (Province of Taranto), the Parco del Marchese water-lifting facility is operated by \acrfull{aqp} and is widely referenced—by  itself and by local press—as the largest pumping station in Europe and the most important node of the Apulian potable-water conveyance system. Functionally, the plant elevates and transfers very high flows toward the central-southern Apulian distribution network; public materials report a nominal hydraulic capacity on the order of $\approx 7,000 ~\dfrac{L}{s}$.\footnote{\href{https://www.aqp.it/aqp-comunica/comunicati-stampa/laterza-rinnovato-e-ampliato-il-fotovoltaico-di-aqp-il-sole-solleva-lacqua-e-migliora-laria}{Rinnovato e ampliato fotovoltaico parco del Marchese}}


Within \acrshort{aqp}’s broader supply architecture, Parco del Marchese serves as the strategic interface for water originating from the Sinni scheme. In particular, the main canal of the Sele–Calore system receives flows conveyed via the Parco del Marchese pumping station—located “in agro di Laterza”—and routed onward through the Gioia del Colle hydraulic node toward the downstream branches that feed much of Apulia\footnote{\href{https://it.wikipedia.org/wiki/Acquedotto_pugliese}{\acrshort{aqp}'s Wikipedia page}}. 

This functional placement is consistent with AQP’s description of the Sinni water-treatment works (also located in the Laterza area) as one of the largest of its type in Europe, underscoring the scale of upstream flows that Parco del Marchese must manage.\footnote{\href{https://www.aqp.it/scopri-acquedotto/gli-impianti/potabilizzatori}{\acrshort{aqp} article about describing Sinni Water-treatment plant in Laterza}} 

From an energy and sustainability standpoint, the site has been progressively equipped with on-site photovoltaic generation. The original 1 MW array, inaugurated in 2010, was renewed and expanded in July 2025 to a total installed capacity of 2.6 MW, with an expected annual output of approximately 3.7 GWh. By the end of 2024, the plant had cumulatively produced more than 18 GWh of “green” electricity; according to coverage of \acrshort{aqp}’s announcement, this on-site production contributes materially to the energy needs of a uniquely large pumping facility of $\approx 7,000 ~\dfrac{L}{s}$ and, during peak-sun hours, can cover up to roughly a quarter of the station’s instantaneous demand. 


Beyond day-to-day operations, \acrshort{aqp} highlights Parco del Marchese in its public-engagement activities (guided technical visits) and routinely manages site upkeep and security through formal procurement, which further indicates the facility’s centrality within the company’s operational portfolio. Examples include AQP’s “Visita gli impianti” program listing the Parco del Marchese lifting station among visitable sites, and recent notices covering extraordinary maintenance of external areas and surveillance services for the associated photovoltaic infrastructure.